\begin{enumerate}[label=\textbf{\ansnum.\arabic*}]
  \item
\begin{verbatim}
def safe_mode_change(client, mode):
    res = client.ChangeMode(mode)
    if res != 0:          # 0 = success in this SDK version (check!)
        raise RuntimeError(f"ChangeMode({mode}) failed with code {res}")
\end{verbatim}
  \item The walking controller accelerates and decelerates (the first and last steps are shorter); the real speed is lower than the command on this floor (friction, the controller's limits); the command was not repeated often enough and the robot paused. Also the tape measurement itself (which point of the foot).
  \item Turning time $\frac{\pi/2}{0.5}=\pi\approx3.14$\,s.
\begin{verbatim}
for side in range(4):
    for v, w, secs in [(0.25, 0, 4.0), (0, 0.5, 3.14)]:
        t0 = time.monotonic()
        while time.monotonic() - t0 < secs:
            client.Move(v, 0.0, w)
            time.sleep(0.1)
client.Move(0, 0, 0)
\end{verbatim}
The real robot does not reach the commanded speeds exactly, it slips, and errors in the angle add up over four turns; there is no feedback of where it is. That is why later chapters use odometry and localisation.
  \item
\begin{verbatim}
def look_at(client, yaw, pitch=0.0, limit=1.0):
    yaw = max(-limit, min(limit, yaw))
    client.RotateHead(pitch, yaw)
\end{verbatim}
  \item The walking controller needs time to switch modes, and a large step command at once can make it fall. Wait until the mode has changed (read it back if the SDK allows), and \emph{ramp} the velocity: increase it in small steps over about a second.
  \item Open-ended; no measurement was available when this book was drafted. Expect a small sway while standing, a larger deflection with a push and an oscillation that dies out; the numbers for your robot are the result. Keep the plot: it is a reference for later changes to the standing controller.
\end{enumerate}
